\documentclass{zkdl-presentation-template}


% Circuit diagrams: inputs gray, gates indigo (+), orange (x), wires gray
\colorlet{circle edge}{oc-gray-7}
\colorlet{circle area}{oc-gray-2}
\colorlet{gate1 edge}{oc-indigo-7}
\colorlet{gate1 area}{oc-indigo-1}
\colorlet{gate2 edge}{oc-orange-7}
\colorlet{gate2 area}{oc-orange-1}
\colorlet{wire}{oc-gray-5}
\colorlet{wire out}{oc-gray-7}
\tikzset{
    var/.style={circle, draw=circle edge, fill=circle area, very thick, minimum size=1cm, text centered},
    gate1/.style={circle, draw=gate1 edge, fill=gate1 area, ultra thick, minimum size=1cm, text centered},
    gate2/.style={circle, draw=gate2 edge, fill=gate2 area, ultra thick, minimum size=1cm, text centered},
    arrow/.style={-Stealth, ultra thick}
}

\title[UltraGroth]{\textbf{UltraGroth. Lookup Checks Enabled in Groth16}}
\author{Blockstream Research}
\date{\today}
\homepage{research.blockstream.com}
\github{BlockstreamResearch/zk-book}

\begin{document}
    \maketitleframe

    \begin{frame}{Plan}
        \tableofcontents
    \end{frame}

	\section{Introduction}

    \begin{frame}{Motivation}
        We typically want to check \textbf{inclusion $\{z_i\}_{i \in [n]} \subseteq
        \{t_j\}_{j \in [v]}$}, where $\{z_i\}_{i \in [n]}$ is part of the witness
        while $\{t_j\}_{j \in [v]}$ is the lookup table.\pause

        \textcolor{oc-green-8}{\textbf{Example usage:}} effective range-checks
        (Bionetta, Rarimo circuits, non-native ZK verifications etc.)


        \begin{tikzpicture}[
            % Style for the bit squares
            bit/.style={draw=oc-gray-9, very thick, minimum size=0.6cm, outer sep=0},
            % Style for the labels below the rows
            info/.style={font=\small}
        ]

        % --- First line: 16 individual bits ---
        \node at (-1, 2.0) {$x=$};
        \foreach \i in {0,...,15} {
            \node[bit] at (\i*0.6, 2.0) {};
            \node at (\i*0.6, 2.0) {$x_{\i}$};
        }
        \node[info] at (7.5*0.6, 1.2) {16 constraints};

        % --- Second line: 4 chunks of 4 bits ---
        \node at (-1, 0) {$x=$};
        \foreach \i/\chunkcolor/\chunktext in {0/oc-indigo-1/oc-indigo-8, 1/oc-red-1/oc-red-8, 2/oc-green-1/oc-green-8, 3/oc-orange-1/oc-orange-8} {
            \fill[\chunkcolor] (\i*4*0.6 - 0.3, -0.3) rectangle (\i*4*0.6 + 4*0.6 - 0.3, 0.3);
            \node[text=\chunktext, font=\bfseries] at (\i*4*0.6 + 1.5*0.6, -0.8) {$x_{\i}$};
        }
        \foreach \i in {0,...,15} {
            \node[bit] at (\i*0.6, 0) {};
        }
        \foreach \i in {1,2,3} {
            \draw[thick, oc-gray-9] (\i*4*0.6 - 0.3, -0.3) -- (\i*4*0.6 - 0.3, 0.3);
        }
        \node[info, text width=10cm, align=center] at (7.5*0.6, -1.5) {\textcolor{oc-green-8}{4 constraints} + \textcolor{oc-red-8}{one-time $2^4$ commitment}};

        \end{tikzpicture}

        \pause\textbf{Generally:} for $n$-bit range-check, the circuit's complexity
        reduces from $O(n)$ to $O(2^w + \frac{n}{w})$, which
        yields $O(n/\log n)$ asymptotic.
    \end{frame}

    \begin{frame}{Logup Check}
        \begin{theorem}[On fractional lookup check]
            The inclusion check $\{z_i\}_{i \in [n]} \subseteq \{t_i\}_{i \in
            [v]}$ is satisfied if and only if there exists the set of
            multiplicities $\{\mu_i\}_{i \in [v]}$ where $\mu_i = \#\{j \in [n]:
            z_j = t_i\}$ such that for $\gamma \leftarrowS \mathbb{F}$:
            \begin{equation*}
                \sum_{i \in [n]} \frac{1}{\gamma + z_i} = \sum_{i \in [v]} \frac{\mu_i}{\gamma + t_i}
            \end{equation*}
        \end{theorem}

        \pause\textit{Naive approach:} Define \texttt{\textcolor{oc-pink-7}{signal}} \texttt{gamma} and implement this check
        in-circuit. This costs exactly $n + 2v$ constraints.\pause

        \textcolor{oc-red-8}{\textbf{Problem:}} We \textit{cannot define
        random} signals in Circom since Groth16, compared to
        $\mathcal{P}$lon$\mathcal{K}$ or sumcheck-based approaches, is not
        compiled from interactive protocol using Fiat-Shamir heuristic.
    \end{frame}

    \begin{frame}{Other implications}
        UltraGroth, though, is not about lookup checks only.\pause

        Assume that you need to implement multiplication of two matrices $A,B
        \in \mathbb{F}^{n \times n}$. \textbf{Naive way} is to compute $C=AB$ by definition:\pause
        \begin{equation*}
            C_{i,j} = \sum_{k \in [n]} A_{i,k}B_{k,j} \hspace{30px} \text{\commentgray{Costs $n$ constraints per $C_{i,j}$}}
        \end{equation*}
        As we have $n^2$ elements in $C$, we thus need $n^3$ constraints.\pause

        We can instead apply the \textbf{Freivalds' protocol}. Sample random
        $\boldsymbol{\gamma} \leftarrowS \mathbb{F}^n$, compute $C$ off-circuit
        and then verify:
        \begin{equation*}
            AB\boldsymbol{\gamma} = C\boldsymbol{\gamma} \hspace{30px} \text{\commentgray{Costs $3n^2$ constraints}}
        \end{equation*}

        \textbf{Example:} Attention layer implementation in zkML.
    \end{frame}

    \begin{frame}{Roadmap}
        \begin{enumerate}
            \item We recap the Groth16 construction.
            \item We identify how to make it interactive.
            \item We specify how Fiat-Shamir transformation should be applied.
            \item We show how it can be practically implemented.
        \end{enumerate}
    \end{frame}

    \section{Groth16 Recap}

    \begin{frame}{R1CS}
        Recall that we can encode any NP statement in the form of $m$ equations
        of form $\langle \boldsymbol{\ell}_j, \mathbf{z} \rangle \cdot \langle
        \boldsymbol{r}_j, \mathbf{z} \rangle = \langle \boldsymbol{o}_j,
        \mathbf{z} \rangle$ for $j \in [m]$ and $\mathbf{z} \in \mathbb{F}^n$.
        Such way of representing the statement is called \textbf{R1CS
        arithmetization}.\pause

        This equality is rewritten more succinctly in the matrix form:
        \begin{equation*}
            L\mathbf{z} \odot R\mathbf{z} = O\mathbf{z}
        \end{equation*}

        \pause\begin{itemize}[label=\cmark]
            \item As of now, this is one of the most optimal arithmetization
            systems available (compared to PlonK and AIR).
            \item All linear operations over elements of $\mathbf{z}$ cost
            \textbf{0 constraints}, compared to PlonK.
        \end{itemize}
    \end{frame}

    \begin{frame}{R1CS: Example}
        \begin{example}
            Suppose the program computes the expression $y=x_1^3 + x_2^2$. \pause

            \vspace{10px}
            \begin{columns}[onlytextwidth]
                \begin{column}{0.45\textwidth}
                    \begin{center}
                        \underline{\textbf{Circuit Diagram}}

                        \scalebox{0.62}{
                        \begin{tikzpicture}
                            % Nodes
                            \node[var] (x1) at (0, 1) {$x_1$};
                            \node[gate2] (x1_x1) at (2, 1) {$\times$};
                            \node[gate2] (x1_x1_x1) at (4, 1) {$\times$};

                            \node[var] (x2) at (0, -1) {$x_2$};
                            \node[gate2] (x2_x2) at (2, -1) {$\times$};

                            \node[gate1] (plus) at (5.0, -0.5) {$+$};

                            % x1**3
                            \draw[arrow,wire] (x1) to [bend left=45] (x1_x1);
                            \draw[arrow,wire] (x1) to [bend right=15] (x1_x1);
                            \draw[arrow,wire] (x1_x1) -- node[midway, above] {\textcolor{oc-blue-8}{$t_1$}} (x1_x1_x1);
                            \draw[arrow,wire] (x1) to [bend right=45] (x1_x1_x1);

                            % x2**2
                            \draw[arrow,wire] (x2) to [bend left=30] (x2_x2);
                            \draw[arrow,wire] (x2) to [bend right=30] (x2_x2);

                            % Summation
                            \draw[arrow,wire] (x1_x1_x1) -- node[midway, right] {\textcolor{oc-green-8}{$t_2$}} (plus);
                            \draw[arrow,wire] (x2_x2) -- node[midway, below] {\textcolor{oc-violet-7}{$t_3$}} (plus);

                            % Result
                            \node[var] (q) at (7.0, -0.5) {$y$};
                            \draw[arrow,wire out] (plus) -- (q);
                        \end{tikzpicture}
                        }
                    \end{center}
                \end{column}

                \begin{column}{0.4\textwidth}
                    \begin{center}
                        \underline{\textbf{Constraints}}
                        \begin{align*}
                            \textcolor{oc-blue-8}{t_1} = x_1 \cdot x_1 \\
                            \textcolor{oc-green-8}{t_2} = \textcolor{oc-blue-8}{t_1} \cdot x_1 \\
                            \textcolor{oc-violet-7}{t_3} = x_2 \cdot x_2 \\
                            y = \textcolor{oc-green-8}{t_2} + \textcolor{oc-violet-7}{t_3}
                        \end{align*}
                    \end{center}
                \end{column}
            \end{columns}

            \vspace{10px}

            \pause In this case, the witness looks as $\mathbf{z} = (1,x_1,x_2,t_1,t_2,t_3,y)$,
            and (for simplicity, consider only $L,R$):
            \begin{gather*}
                \footnotesize L = \begin{bmatrix}
                    0 & 1 & 0 & 0 & 0 & 0 & 0\\
                    0 & 0 & 0 & 1 & 0 & 0 & 0 \\
                    0 & 0 & 1 & 0 & 0 & 0 & 0 \\
                    0 & 0 & 0 & 0 & 1 & 1 & 0
                \end{bmatrix}, \; R = \begin{bmatrix}
                    0 & 1 & 0 & 0 & 0 & 0 & 0\\
                    0 & 1 & 0 & 0 & 0 & 0 & 0 \\
                    0 & 0 & 1 & 0 & 0 & 0 & 0 \\
                    1 & 0 & 0 & 0 & 0 & 0 & 0
                \end{bmatrix}
            \end{gather*}
        \end{example}
    \end{frame}

    \begin{frame}{QAP}
        We interpolate the columns of each of the matrices, thus getting $3n$
        polynomials $\{(\ell_i(T),r_i(T),o_i(T))\}_{i \in [n]} \subseteq
        \mathbb{F}[T]^{\leq m}$:
        \begin{equation*}
            \ell_i(\omega^j) = L_{j,i}, \quad r_i(\omega^j) = R_{j,i}, \quad o_i(\omega^j) = O_{j,i}, \quad i \in [n], \; j \in [m]
        \end{equation*}

        \pause Now, the same R1CS check can be encoded over polynomial space:
        \begin{equation*}
            \sum_{i \in [n]}z_i\ell_i(T) \cdot \sum_{i \in [n]} z_ir_i(T) = \sum_{i \in [n]}z_io_i(T) + t_{\Omega}(T)h(T),
        \end{equation*}
        where $h(T)$ is computed by a prover and $t_{\Omega}(T) \triangleq
        \prod_{h \in \Omega}(T-h)$ is the \textit{vanishing polynomial} over
        \textit{evaluation domain $\Omega = \{\omega^j\}_{j \in [m]}$}. The
        corresponding relation:\pause
        \begin{equation*}
            \footnotesize\mathcal{R}_{\text{QAP}} = \left\{
        \begin{aligned}
            &\mathbbm{x} = \{z_i\}_{i \in \mathcal{I}_{X}} \\
            &\mathbbm{w} = \{z_i\}_{i \in \mathcal{I}_{W}}
        \end{aligned}
        \; \middle| \;
        \begin{matrix}
            \sum_{i \in [n]}z_i\ell_i(T) \cdot \sum_{i \in [n]}z_ir_i(T) = \sum_{i \in [n]}z_io_i(T) + t_{\Omega}(T)h(T) \\
            \text{for some} \; h(T) \in \mathbb{F}[T]
        \end{matrix}
    \right\}
        \end{equation*}
    \end{frame}

    \begin{frame}{Linear non-interactive proofs}
        Recall that Groth16 is compiled from \textit{Linear non-interactive proofs}.\pause
        \begin{definition}[Linear Non-Interactive Proof]
            The \textbf{Linear Non-Interactive Proof} consists of the following procedures:
            \begin{itemize}
                \item $\textsf{Setup}(1^{\lambda}, \mathcal{R}) \to (\boldsymbol{\sigma},\boldsymbol{\tau})$. The setup
                returns $\boldsymbol{\sigma} \in \mathbb{F}^m$ and $\boldsymbol{\tau} \in \mathbb{F}^n$. \pause
                \item $\textsf{Prove}(\boldsymbol{\sigma}, \mathbbm{x},
                \mathbbm{w}) \to \boldsymbol{\pi}$. $\mathcal{P}$ chooses the
                matrix $\Pi \in \mathbb{F}^{k \times m}$ and computes the proof
                as $\boldsymbol{\pi} \gets \Pi \boldsymbol{\sigma}$.\pause
                \item $\textsf{Verify}(\boldsymbol{\sigma}, \mathbbm{x}, \pi)\to\{0,1\}$.
                The verifier gets the arithmetic circuit $t: \mathbb{F}^{m+k}
                \to \mathbb{F}^{\eta}$ of degree $d$ and verifies whether
                $t(\boldsymbol{\sigma},\boldsymbol{\pi}) = 0$.\pause
            \end{itemize}
        \end{definition}
        \vspace{-10px}

        \textbf{Groth16} is essentially a Linear NIP where $d=2$ and
        $\boldsymbol{\sigma}$ is given by:
        \begin{equation*}
            \boldsymbol{\sigma} = \left(\alpha,\beta,\gamma,\delta,\{\tau^i\}_{i \in [n]}, \left\{\frac{f_i(\tau)}{\gamma}\right\}_{i \in [m]}, \left\{\frac{\tau^it_{\Omega}(\tau)}{\delta}\right\}_{i \in [n]}\right),
        \end{equation*}
        with $f_i(T) := \beta \ell_i(T) + \alpha r_i(T) + o_i(T)$ and
        $\boldsymbol{\tau} = (\alpha,\beta,\gamma,\delta,\tau)$.
    \end{frame}

    \begin{frame}{Groth16 Construction}
        Fix bilinear group $\mathcal{G} = (\mathbb{G}_1, \mathbb{G}_2,
        \mathbb{G}_T,e)$ with pairing $e: \mathbb{G}_1 \times \mathbb{G}_2 \to
        \mathbb{G}_T$.\pause

        \begin{itemize}
            \item $\mathsf{Setup}(1^{\lambda},\mathcal{R}_{\text{QAP}}) \to (\mathsf{pp},\mathsf{vp})$. See previous slide.\pause
            \item $\mathsf{Prove}(\mathsf{pp},\mathbbm{x}, \mathbbm{w}) \to \boldsymbol{\pi}$. Sample random
            $r,s \leftarrowS \mathbb{F}$ and output $\boldsymbol{\pi} \gets (g_1^{a(\tau)}, g_1^{c(\tau)}, g_2^{b(\tau)})$ where:
            \begin{gather*}
                a(T) = \alpha + \sum_{i \in [n]}z_i\ell_i(T) + r\delta, \; b(T) = \beta + \sum_{i \in [n]}z_ir_i(T) + s\delta, \\
                c(T) = \delta^{-1}\left(\sum_{i \in \mathcal{I}_W} z_if_i(T)+h(T)t_{\Omega}(T)\right) + a(T)s + b(T)r - rs\delta
            \end{gather*}
            \pause \item $\mathsf{Verify}(\mathsf{vp},\mathbbm{x},\pi) \to \{0,1\}$. Parse $\pi=(\pi_A,\pi_C,\pi_B)$ and
            accept the proof if and only if
            \begin{equation*}
                e(\pi_A,\pi_B) = e(g_1^{\alpha},g_2^{\beta}) \cdot e(g_1^{\iota(\tau)}, g_2^{\gamma}) \cdot e(\pi_C, g_2^{\delta}),
            \end{equation*}
            where $\iota(T) := \gamma^{-1}\sum_{i \in
            \mathcal{I}_X}z_if_i(T)$ is the input commitment.
        \end{itemize}
    \end{frame}

    \section{UltraGroth}

    \begin{frame}{Desired Interactive Protocol}
        We would like to have the following interactive protocol (IP) between
        the prover $\mathcal{P}$ and verifier $\mathcal{V}$.\pause
        \begin{block}{Interactive Protocol}
        \underline{\textbf{Input:}} Relation $\mathcal{R}_{\text{QAP}}$ and public statement
        $\mathbbm{x}_0$.\pause

        \underline{\textbf{Round 0:}} $\mathcal{P}$ runs the circuit without imposing lookup
        check and gets witness $\mathbbm{w}_0$. $\mathcal{V}$ sends the random
        challenge $\mathbbm{x}_1 \leftarrowS \mathbb{F}$.\pause

        \underline{\textbf{Round 1:}} $\mathcal{P}$ computes the second part of the witness
        $\mathbbm{w}_1$, corresponding to the lookup check $\sum_{i \in
        [n]}\frac{1}{\mathbbm{x}_1+z_i} = \sum_{i \in
        [v]}\frac{\mu_i}{\mathbbm{x}_1+t_i}$. The prover $\mathcal{P}$ sends
        $\mathbbm{w}_1$ and $h(T)$ to verifier.\pause

        \underline{\textbf{Check:}} $\mathcal{V}$ checks $\ell(T)r(T) = o(T) +
        t_{\Omega}(T)h(T)$.\pause
        \end{block}

        \textbf{Compiling IP into NIZK.} Apply Fiat-Shamir transformation:
        sample challenge as $\mathbbm{x}_1=\mathsf{H}(\mathsf{pp},
        \mathbbm{x}_0, \mathbbm{w}_0)$.\pause

        \textcolor{oc-red-8}{\textbf{Problem.}} We cannot practically
        ``hash'' the witness part $\mathbbm{w}_0$.
    \end{frame}

    \begin{frame}{2-round UltraGroth Construction}
        Split public indexing set $\mathcal{I}_X$ into two parts:
        $\mathcal{I}_X^{\langle 0 \rangle}$ and $\mathcal{I}_X^{\langle 1
        \rangle}$. Similarly, split the witness indexing set $\mathcal{I}_W$
        into $\mathcal{I}_W^{\langle 0 \rangle}$ and $\mathcal{I}_W^{\langle 1
        \rangle}$.\pause

        \underline{\textbf{Input:}} Relation $\mathcal{R}_{\text{QAP}}$ and public statement
        $\mathbbm{x}_0$.\pause

        \underline{\textbf{Round 0:}} $\mathcal{P}$ runs circuit without lookup
        check and gets witness $\mathbbm{w}_0$. She samples $r_0 \leftarrowS
        \mathbb{F}$, and computes $\pi_C^{\langle 0 \rangle} \gets
        g_1^{c_0(\tau)}$ as:
        \begin{equation*}
            c_0(T) = \delta_0^{-1}\sum_{j \in \mathcal{I}_W^{\langle 0 \rangle}} z_jf_j(T) + r_0\delta
        \end{equation*}

        \pause\underline{\textbf{Round 1:}} $\mathcal{P}$ samples the challenge $\mathbbm{x}_1 \gets \mathsf{H}(\mathsf{pp},\pi_C^{\langle 0 \rangle})$,
        samples $r,s \leftarrowS \mathbb{F}$ and computes $\pi_C^{\langle 1 \rangle} \gets g_1^{c_1(\tau)}$ as:
        \begin{equation*}
            c_1(T) = \delta^{-1}\left(\sum_{j \in \mathcal{I}_W^{\langle 1 \rangle}} z_jf_j(T)+h(T)t_{\Omega}(T)\right) + a(T)s + b(T)r - r_0\delta_0 - rs\delta
        \end{equation*}
    \end{frame}

    \begin{frame}{2-round UltraGroth Construction: The rest}
        Then, parts $\pi_A \gets g_1^{a(\tau)}$ and $\pi_B \gets g_2^{b(\tau)}$ are computed as usual via:
        \begin{equation*}
            a(T) = \alpha + \sum_{i \in [n]}z_i\ell_i(T) + r\delta, \; b(T) = \beta + \sum_{i \in [n]}z_ir_i(T) + s\delta.
        \end{equation*}

        \pause\textbf{Note:} $\delta_0c_0(T) + \delta c_1(T)$ is exactly $\delta c(T)$ in the original Groth16. \pause
        Thus, $\mathcal{V}$ checks:
        \begin{equation*}
            e(\pi_A,\pi_B) = e(g_1^{\alpha},g_2^{\beta}) \cdot e(g_1^{\iota(\tau)},g_2^{\gamma}) \cdot e(\pi_C^{\langle 0 \rangle}, g_2^{\delta_0}) \cdot e(\pi_C^{\langle 1 \rangle}, g_2^{\delta}),
        \end{equation*}
        where $\iota(T) = \gamma^{-1}\sum_{i \in \mathcal{I}_X}z_if_i(T)$ as before
        and $\mathbbm{x}_1 = \mathsf{H}(\mathsf{pp}, \pi_C^{\langle 0
        \rangle})$.

        \pause\begin{remarkblock}{Conclusion}
            UltraGroth protocol's verifier is only \textbf{4 pairings},
            \textbf{1 hashing operation}, and $O(|\mathbbm{x}|)$
            exponentiations over $\mathbb{G}_1$.
        \end{remarkblock}
    \end{frame}

    \begin{frame}{Multi-round UltraGroth}
        \begin{definition}[dQAP]
            We define the \textbf{$(d+1)$-round quadratic arithmetic program}
            (or \textbf{$d$QAP}, for short), as follows:
            \begin{equation*}
                \small\mathcal{R}_{\text{dQAP}} = \left\{
            \begin{matrix}
            \mathbbm{x}_i = \{z_j\}_{j \in \mathcal{I}_{X}^{\langle i \rangle}} \\
            \mathbbm{w}_i = \{z_j\}_{j \in \mathcal{I}_{W}^{\langle i \rangle}} \\
            \text{for }i \in [d+1]
        \end{matrix}
        \; \middle| \;
        \begin{matrix}
            \ell(T) \cdot r(T) = o(T) + t_{\Omega}(T)h(T) \\
            \ell(T) = \sum_{i \in [n]}z_i\ell_i(T), \\
            r(T) = \sum_{i \in [n]}z_ir_i(T), \\
            o(T) = \sum_{i \in [n]}z_io_i(T), \\
            \text{for some} \; h(T) \in \mathbb{F}[T]
        \end{matrix}
    \right\},
            \end{equation*}
        where $\{\mathcal{I}_X^{\langle i \rangle}\}_{i \in [d+1]}$ and
        $\{\mathcal{I}_W^{\langle i \rangle}\}_{i \in [d+1]}$ partition $[n]$.\pause
        \end{definition}

        \begin{figure}
            \centering
            \includegraphics[width=1.0\textwidth]{images/lecture_18/wtns_structure.png}
        \end{figure}
    \end{frame}

    \begin{frame}{Strategy}
        \begin{definition}[Strategy]
            Define \textbf{strategy} for $\mathcal{R}_{\text{dQAP}}$ as the
            collection of functions $S=\{S_i\}_{i \in [d]}$ each of which
            computes the witness for the given round given previous witnesses
            and challenges and the current challenge, sampled by the verifier.
            In other words,
            \begin{equation*}
                \mathbbm{w}_i = S_i(\mathbbm{x}_0,\dots,\mathbbm{x}_i,\mathbbm{w}_0,\dots,\mathbbm{w}_{i-1})
            \end{equation*}
        \end{definition}

        \pause\begin{example}
            0QAP represents the regular QAP with the strategy $S = \{S_0\}$ that
            consists of the witness generator: $\mathbbm{w} = S_0(\mathbbm{x})$.\pause
            In turn, 1QAP represents the lookup Groth16 version where
            $\mathbbm{w}_0=S_0(\mathbbm{x}_0)$ computes the witness without
            lookups while
            $\mathbbm{w}_1=S_1(\mathbbm{x}_0,\mathbbm{x}_1,\mathbbm{w}_0)$
            computes lookup constraints.
        \end{example}
    \end{frame}

    \begin{frame}{$d$-round UltraGroth}
        Initialize accumulator $a_0 := \mathsf{H}(\mathsf{pp})$.\pause

        \underline{\textbf{On each round $i \in [d]$:}}
        \begin{itemize}
            \item Sample $r_i \leftarrowS \mathbb{F}$.\pause
            \item Compute witness $\mathbbm{w}_i \gets S_i(\mathbbm{x}_0,\dots,\mathbbm{x}_i,\mathbbm{w}_0,\dots,\mathbbm{w}_{i-1})$.\pause
            \item Compute $\pi_C^{\langle i \rangle}$ with $c_i(T) :=
            \delta_i^{-1}\sum_{j \in \mathcal{I}_W^{\langle i \rangle}}
            z_jf_j(T) + r_i\delta_d$.\pause
            \item Update accumulator $a_{i+1} \gets \mathsf{H}(a_i, \pi_C^{\langle i \rangle})$.\pause
            \item If $i<d$, for each $j \in \mathcal{I}_X^{\langle i+1 \rangle}$,
            set $z_j \gets \mathsf{H}(a_{i+1},g_1^j)$.
        \end{itemize}
    \end{frame}

    \begin{frame}{$d$-round UltraGroth: Last Round}
        \underline{\textbf{During the last round:}}
        \begin{itemize}
            \item Compute $h(T)$ similar to Groth16.\pause
            \item Sample $r,s \leftarrowS \mathbb{F}$ and compute $\pi_A \gets
            g_1^{a(\tau)}$, $\pi_B \gets g_2^{b(\tau)}$, and the last proof
            piece $\pi_C^{\langle d \rangle} \gets g_1^{c_d(\tau)}$ where:
            {\small\begin{gather*}
                a(T) = \alpha + \sum_{i \in [n]} z_i\ell_i(T) + r\delta_d, \; b(T) = \beta + \sum_{i \in [n]}z_ir_i(T) + s\delta_d, \\
                c_d(T) = \delta_d^{-1}\left(\sum_{i \in \mathcal{I}_W^{\langle d \rangle}}z_if_i(T) + h(T)t_{\Omega}(T)\right) + a(T)s+b(T)r-\sum_{i \in [d]}r_i\delta_i - rs\delta_d
            \end{gather*}}
            \pause\item \textbf{Output} proof $\pi = (\pi_A,\pi_B,\{\pi_C^{\langle i
            \rangle}\}_{i \in [d+1]}) \in \mathbb{G}_1 \times \mathbb{G}_2
            \times \mathbb{G}_1^{d+1}$. \pause
        \end{itemize}

        \underline{\textbf{Verification:}} $e(\pi_A,\pi_B) = e(g_1^{\alpha},g_2^{\beta}) \cdot e(g_1^{\iota(\tau)},g_2^{\gamma}) \cdot \prod_{i \in [d+1]}e(\pi_C^{\langle i \rangle}, g_2^{\delta_i})$.
    \end{frame}

    \begin{frame}{UltraGroth Efficiency}
    \textcolor{oc-green-8}{\textbf{Groth16} performance} over
    the circuit of size $n$ and statement size $\ell$.
    \begin{itemize}
        \item \textbf{Prover work}: MSM of size $O(n)$ over
        $\mathbb{G}_1$ and $\mathbb{G}_2$.
        \item \textbf{Proof size:} $2\mathbb{G}_1+\mathbb{G}_2$.
        \item \textbf{Verifier work:} 3 pairings + $O(\ell)$
        $\mathbb{G}_1$ exps.
    \end{itemize}

    \pause\textcolor{oc-blue-8}{\textbf{UltraGroth} performance over $\mathcal{R}_{\text{dQAP}}$} in turn:
    \begin{itemize}
        \item \textbf{Prover work}: MSM of size
        $O(\textcolor{oc-blue-8}{n/\log n})$ over $\mathbb{G}_1$
        and $\mathbb{G}_2$.
        \item \textbf{Proof size:}
        $\textcolor{oc-blue-8}{(d+2)}\mathbb{G}_1+\mathbb{G}_2$.
        \item \textbf{Verifier work:} \textcolor{oc-blue-8}{$(d+3)$}
        pairings + $O(\ell)$ $\mathbb{G}_1$ exps +
        \textcolor{oc-blue-8}{$\sum_{i \in [d+1]\setminus
        \{0\}}|\mathbbm{x}_i|$ hashing operations}.
        \item Allowed interactiveness for potentially more complex protocols.
    \end{itemize}
    \end{frame}

    \makeclosingframe
\end{document}
